\begin{lemma}[Matching count in a regular bipartite graph]
Let $F$ be a finite bipartite graph with disjoint parts $U,W$ of size $n$.
If every vertex on both sides has degree $k$ across the two parts, then $F$
has at least
\[
 (k/n)^n n!
\]
perfect matchings.  The assertion uses the conventions $0^0=1$ and
$0!=1$, so it also includes the empty graph on two empty parts.
\end{lemma}
\begin{proof}
If $n=0$, both parts are empty. There is one function $U\to W$,
it is bijective, and all edge conditions are vacuous. Thus by
\noderef{PerfectMatching} there is exactly one perfect matching.
In extended nonnegative reals $(k/0)^0=1$, independently of $k$,
and $0!=1$, proving the assertion. Suppose $n>0$.
If $k=0$, the lower bound is $0^n n!=0$, so the assertion is immediate.
Henceforth $k>0$ too.

Choose bijections $e_U:U\to\mathrm{Fin}(n)$ and
$e_W:W\to\mathrm{Fin}(n)$, possible because both finite sets have size $n$.
Let $A_{ij}$ be the indicator that $e_U^{-1}(i)e_W^{-1}(j)$ is an edge.
For fixed $i$, summing the indicators over $j$ counts precisely the
neighbors in $W$ of $e_U^{-1}(i)$: the indexing bijection visits each vertex
of $W$ once. The degree hypothesis makes this sum $k$. The analogous
argument on columns uses the right-degree hypothesis and also gives $k$.
All entries of $A$ are nonnegative. By
\noderef{AdjacencyPermanentMatchingCount}, $\operatorname{per}(A)$ is
the natural matching cardinality $M$, cast to $\mathbb R$.
This cited identity uses Mathlib's permutation orientation, so no choice
of row-versus-column convention is being suppressed.

By \noderef{RegularMatrixNormalization}, $B=A/k$ is nonnegative with
all row and column sums one, and
$\operatorname{per}(A)=k^n\operatorname{per}(B)$.
Apply the checked imported theorem
\texttt{EgorychevFalikmanFormalization.VanDerWaerdenPermanentInput}.
It states that for any real $n\times n$ matrix $B$, the hypotheses
$B_{ij}\ge0$ for every $i,j$, $\sum_jB_{ij}=1$ for every $i$, and
$\sum_iB_{ij}=1$ for every $j$ imply
\[
 \operatorname{per}(B)\ge\frac{n!}{n^n}.
\]
All three hypotheses are exactly the conclusions of
\noderef{RegularMatrixNormalization} established above for $B=A/k$.
Thus this is an application of the proved permanent bound, not an appeal
to a bibliographic citation.
Multiply by $k^n>0$ and use the normalization and count identities:
\[
 M=\operatorname{per}(A)\ge k^n\frac{n!}{n^n}
   =(k/n)^n n!.
\]
This last calculation is in $\mathbb R$, with positive denominator $n^n$.
Both $M$ and the lower bound are nonnegative and finite. Apply the monotone
map $\mathrm{ofReal}:\mathbb R\to\mathrm{ENNReal}$. For nonnegative real
arguments it preserves finite products and powers; division by the
positive real integer $n$ becomes division by its finite nonzero
$\mathrm{ENNReal}$ cast. The natural casts of $k,n,n!$, and $M$ agree with
their real casts followed by $\mathrm{ofReal}$. Consequently the resulting
inequality is exactly the stated one, not merely a real-valued variant.
\end{proof}
